\chapter{الجهازان التناسليان}\label{ch:g8:reproductive-systems}

احتاج \Cref{ch:g8:sexual-reproduction} إلى \omterm{def:g8:sexual-reproduction:gametes}{الأمشاج}؛
وشغّل \cref{ch:g8:puberty} مصانعها. ويرسم هذا الفصل
خريطة المصانع نفسها: الجهاز التناسلي الذكري والجهاز
التناسلي الأنثوي --- الأعضاء، \omterm{def:g8:sexual-reproduction:gametes}{والمشيج} الذي ينتجه كل
منهما، والإيقاع الشهري اللافت للجهاز الأنثوي. تشريح
بسيط بأسماء بسيطة: هذه أعضاء \omterm{def:g4:heart-and-blood:heart}{كالقلب} \omterm{def:g7:how-animals-breathe:four}{والرئتين}، وهذا
فصلها.

\section{الجهاز الذكري}

\begin{definition}[الجهاز التناسلي الذكري]\label{def:g8:reproductive-systems:male}
أعضاء الجهاز الذكري:
\begin{itemize}
  \item \emph{الخصيتان}\index{خصية}، محفوظتان خارج دفء
        الجسم في \emph{كيس الصفن} --- وهما مصنع \omterm{def:g8:sexual-reproduction:gametes}{الأمشاج}:
        تنتجان منذ \omterm{def:g8:puberty:puberty}{البلوغ} \emph{النطف}\index{نطفة} إنتاجًا
        مستمرًا، بمئات الملايين؛
  \item قنوات تخزّن النطف وتنقلها، تنضم إليها في الطريق
        غدد تغذّي سوائلها النطف وتحملها --- فتتكوّن معًا
        \emph{المني}؛
  \item \emph{القضيب}\index{قضيب}، ومنه يخرج المني ---
        وهو نفسه العضو الذي يحمل \omterm{def:g7:removing-wastes:kidneys}{البول} في مهمته الأخرى
        (\cref{def:g7:removing-wastes:kidneys})؛ ويفصل
        ترتيب من الصمامات بين المهمتين فصلًا تامًا.
\end{itemize}
\end{definition}

\begin{example}[النطفة عن قرب]\label{ex:g8:reproductive-systems:sperm}
تحت \omterm{def:g5:first-look-at-cells:microscope}{المجهر}
(\cref{def:g5:first-look-at-cells:microscope})، تبدو
\omterm{def:g8:reproductive-systems:male}{النطفة} \omterm{def:g5:first-look-at-cells:cell}{خلية} جُرّدت من كل شيء إلا السفر: رأس يكاد يكون
نواةً كلها --- نصيب الأب، مركّزًا مضغوطًا
(\cref{rem:g5:first-look-at-cells:nucleus}) --- وذيل
يضرب الماء. أصغر \omterm{def:g5:first-look-at-cells:cell}{خلايا} الجسم، مبنية للسباحة، مزوّدة بزاد
أيام. ويمتد الإنتاج من \omterm{def:g8:puberty:puberty}{البلوغ} إلى \omterm{def:g3:stages-of-life:stages}{الشيخوخة}: فمصنع \omterm{def:g4:animal-reproduction:sexual}{الذكر}
مستمر.
\end{example}

\section{الجهاز الأنثوي}

\begin{definition}[الجهاز التناسلي الأنثوي]\label{def:g8:reproductive-systems:female}
أعضاء الجهاز الأنثوي، وكلها في أسفل البطن:
\begin{itemize}
  \item \emph{المبيضان}\index{مبيض} --- مخزن \omterm{def:g8:sexual-reproduction:gametes}{الأمشاج}
        ومصنعها: خلافًا \omterm{def:g8:reproductive-systems:male}{للخصيتين}، يحملان مخزونهما من
        \emph{\omterm{def:g4:animal-reproduction:sexual}{البويضات}} المقبلة منذ \omterm{def:g2:born-grow-die:cycle}{الولادة}، ويطلقانه منذ
        \omterm{def:g8:puberty:puberty}{البلوغ} واحدة واحدة؛
  \item \emph{قناتا البيض}\index{قناة البيض}، أنبوبان
        مهدّبان يلتقطان \omterm{def:g4:animal-reproduction:sexual}{البويضة} المنطلقة --- وهما ملتقى
        الخطوة 2 في \cref{prop:g8:sexual-reproduction:core}؛
  \item \emph{\omterm{prop:g5:human-reproduction:plan}{الرحم}} (\omterm{prop:g5:human-reproduction:plan}{رحم} \cref{prop:g5:human-reproduction:plan})،
        وجداره العضلي السميك وبطانته المتجددة هما مسكن
        \omterm{prop:g8:sexual-reproduction:core}{اللاقحة}؛
  \item \emph{المهبل}، يصل \omterm{prop:g5:human-reproduction:plan}{الرحم} بالخارج --- وفيه يُستقبل
        المني، وهو قناة \omterm{def:g2:born-grow-die:cycle}{الولادة} في
        \cref{ex:g5:human-reproduction:birth}.
\end{itemize}
\end{definition}

\begin{example}[البويضة عن قرب]\label{ex:g8:reproductive-systems:egg}
\omterm{def:g4:animal-reproduction:sexual}{البويضة} نقيض \omterm{def:g8:reproductive-systems:male}{النطفة} في كل شيء إلا الدور: فهي
\emph{أكبر} \omterm{def:g5:first-look-at-cells:cell}{خلايا} الجسم --- عُشر مليمتر، على حافة ما
تراه \omterm{def:g1:the-five-senses:sense}{العين} --- مستديرة، ساكنة، مزوّدة بزاد أيام الرحلة
الأولى (مبدأ المح في \cref{rem:g2:animals-grow-up:egg}،
على مقياس \omterm{def:g5:first-look-at-cells:cell}{الخلية}). مصنع يجري بالملايين بلا انقطاع؛
والآخر يطلق \omterm{def:g5:first-look-at-cells:cell}{خلية} واحدة مزوّدة في الشهر: إنها موازنة
العدد والرعاية في
\cref{prop:g4:animal-reproduction:strategies}، مكتوبة في
\omterm{def:g8:sexual-reproduction:gametes}{الأمشاج} نفسها.
\end{example}

\section{الدورة}

\begin{proposition}[الدورة الشهرية]\label{prop:g8:reproductive-systems:cycle}
من \omterm{def:g8:puberty:puberty}{البلوغ} إلى نحو الخمسين، يُجري الجهاز الأنثوي برنامجًا
متكررًا مدته نحو \emph{28 يومًا} --- هي \emph{الدورة
الشهرية}\index{دورة شهرية}، تنظمها سلسلة \omterm{def:g8:puberty:hormones}{الهرمونات} في
\cref{def:g8:puberty:hormones}:
\begin{enumerate}
  \item تعيد بطانة \omterm{prop:g5:human-reproduction:plan}{الرحم} بناء نفسها، سميكة غنية \omterm{prop:g4:journey-of-food:absorb}{بالدم} ---
        فالمسكن مهيأ؛
  \item في منتصف الدورة، نحو اليوم 14، يطلق أحد \omterm{def:g8:reproductive-systems:female}{المبيضين}
        \omterm{def:g4:animal-reproduction:sexual}{بويضة} واحدة --- هي \emph{الإباضة}\index{إباضة}؛
        فتلتقطها \omterm{def:g8:reproductive-systems:female}{قناة البيض}، وتظل نحو يوم واحد قابلة
        \omterm{def:g4:animal-reproduction:sexual}{للإخصاب}؛
  \item إن لم يحدث \omterm{def:g4:animal-reproduction:sexual}{إخصاب}: فالبطانة غير لازمة؛ تتفكك وتخرج
        عبر المهبل في أيام قليلة --- وهو \emph{الحيض}\index{حيض}
        --- ثم يبدأ البرنامج من جديد؛
  \item وإن حدث \omterm{def:g4:animal-reproduction:sexual}{إخصاب}: تستقر \omterm{prop:g8:sexual-reproduction:core}{اللاقحة} في البطانة المهيأة
        (\cref{def:g5:human-reproduction:pregnancy})،
        فتتوقف الدورة، ويبدأ \omterm{def:g5:human-reproduction:pregnancy}{الحمل}.
\end{enumerate}
\end{proposition}
\admitted

\begin{omfigure}
\begin{tikzpicture}[scale=0.94]
  % cycle wheel
  \draw[very thick, omDef!60] (0,0) circle (2.3);
  % day marks
  \foreach \a/\d in {90/1, 0/8, -90/14, 180/22} {
    \fill[omDef] (\a:2.3) circle (0.07);
  }
  \node[font=\footnotesize, above] at (90:2.45) {اليوم 1};
  \node[font=\footnotesize, right] at (0:2.45) {اليوم 8};
  \node[font=\footnotesize, below] at (-90:2.45) {اليوم 14};
  \node[font=\footnotesize, left] at (180:2.45) {اليوم 22};
  % phases
  \draw[very thick, omThm] (90:2.3) arc[start angle=90, end angle=35,
    radius=2.3];
  \node[font=\footnotesize, omThm, align=center] at (58:3.5)
    {الحيض:\\ تتساقط البطانة};
  \node[font=\footnotesize, omMeth, align=center] at (-35:4.0)
    {إعادة بناء البطانة،\\ سميكة جاهزة};
  \fill[omProp] (-90:2.3) circle (0.12);
  \node[font=\footnotesize, omProp, align=center] at (-90:3.55)
    {الإباضة:\\ بويضة واحدة تنطلق};
  % center
  \node[font=\small, align=center] at (0,0.25)
    {الدورة:\\ نحو 28 يومًا};
  \node[font=\footnotesize, align=center] at (0,-0.75)
    {(الإخصاب يوقفها:\\ الحمل)};
  % arrow direction
  \draw[->, thick, omDef] (150:2.3) arc[start angle=150,
    end angle=120, radius=2.3];
\end{tikzpicture}

{\small \omterm{prop:g8:reproductive-systems:cycle}{الدورة الشهرية} كعجلة: يفتتحها \omterm{prop:g8:reproductive-systems:cycle}{الحيض}، ثم تُبنى
البطانة، وتتوّجها \omterm{prop:g8:reproductive-systems:cycle}{الإباضة} في منتصف الدورة --- \omterm{def:g5:flower-to-fruit:steps}{والإخصاب}
هو الحدث الوحيد الذي يوقف العجلة.}
\end{omfigure}

\begin{remark}[العيش مع الدورة]\label{rem:g8:reproductive-systems:living}
أمور عملية، تقال بصراحة: الدورات غير منتظمة في السنوات
الأولى --- فالآلة تضبط نفسها؛ ويختلف طولها من شخص إلى
آخر ومن شهر إلى آخر؛ وقد يصحب \omterm{prop:g8:reproductive-systems:cycle}{الحيض} تقلصات (\omterm{prop:g5:human-reproduction:plan}{فالرحم} \omterm{def:g3:bones-and-muscles:muscle}{عضلة}،
\cref{def:g3:bones-and-muscles:muscle}، وهو يعمل أثناء
التساقط) --- والدفء والحركة ونصيحة ممرضة المدرسة عند
الحاجة كلها تساعد؛ والفوط والسدادات والكؤوس ليست إلا
العدة العملية للحيض، وهي رف عادي في كل صيدلية. وليس شيء
من هذا مرضًا: فالدورة هي الجهاز يعمل تمامًا كما بُني.
\end{remark}

\begin{method}[قراءة الجهازين جنبًا إلى جنب]\label{met:g8:reproductive-systems:sides}
لتبقى الخريطة واضحة، قارن حسب الوظيفة:
\begin{enumerate}
  \item مصنع \omterm{def:g8:sexual-reproduction:gametes}{الأمشاج}: \omterm{def:g8:reproductive-systems:male}{الخصيتان} --- \omterm{def:g8:reproductive-systems:female}{المبيضان}؛
  \item \omterm{def:g8:sexual-reproduction:gametes}{المشيج}: ملايين سابحة بلا انقطاع --- \omterm{def:g4:animal-reproduction:sexual}{بويضة} واحدة
        مزوّدة في الشهر؛
  \item مكان اللقاء: \omterm{def:g8:reproductive-systems:female}{قناة البيض}، تبلغها \omterm{def:g8:reproductive-systems:male}{النطفة} سباحةً من
        المهبل عبر \omterm{prop:g5:human-reproduction:plan}{الرحم}؛
  \item المسكن: لا مسكن في الخريطة الذكرية --- \omterm{prop:g5:human-reproduction:plan}{والرحم}،
        المهيأ شهريًا، في الخريطة الأنثوية؛
  \item الجدول الزمني: مستمر منذ \omterm{def:g8:puberty:puberty}{البلوغ} --- ودوري، من
        \omterm{def:g8:puberty:puberty}{البلوغ} إلى نحو الخمسين.
\end{enumerate}
\end{method}

\section{تمارين}

\begin{exercise}[$\star$]\label{exo:g8:reproductive-systems:1}
ارسم خريطة الجهاز الذكري: سمّ كل عضو ووظيفته.
\end{exercise}

\begin{exercise}[$\star$]\label{exo:g8:reproductive-systems:2}
ارسم خريطة الجهاز الأنثوي كذلك.
\end{exercise}

\begin{exercise}[$\star$]\label{exo:g8:reproductive-systems:3}
قابل بين \omterm{def:g8:reproductive-systems:male}{النطفة} \omterm{def:g4:animal-reproduction:sexual}{والبويضة}: الحجم، والبنية، والعدد،
والتزويد بالزاد.
\end{exercise}

\begin{exercise}[$\star$]\label{exo:g8:reproductive-systems:4}
ماذا يحدث عند \omterm{prop:g8:reproductive-systems:cycle}{الإباضة}، وفي أي يوم تقريبًا، وكم تبقى
\omterm{def:g4:animal-reproduction:sexual}{البويضة} قابلة \omterm{def:g4:animal-reproduction:sexual}{للإخصاب}؟
\end{exercise}

\begin{exercise}[$\star$]\label{exo:g8:reproductive-systems:5}
ما \omterm{prop:g8:reproductive-systems:cycle}{الحيض}، بلغة البطانة --- وماذا يقول وصوله عن تلك
الدورة؟
\end{exercise}

\begin{exercise}[$\star$]\label{exo:g8:reproductive-systems:6}
أي حدث يوقف الدورة، وماذا يتبعه في \omterm{prop:g5:human-reproduction:plan}{الرحم}؟
\end{exercise}

\begin{exercise}[$\star\star$]\label{exo:g8:reproductive-systems:7}
أين تحدث الخطوات الثلاث في
\cref{prop:g8:sexual-reproduction:core} على خريطة هذا
الفصل --- عضوًا عضوًا؟
\end{exercise}

\begin{exercise}[$\star\star$]\label{exo:g8:reproductive-systems:8}
اقرأ المشيجين بوصفهما الاستراتيجيتين في
\cref{prop:g4:animal-reproduction:strategies} مصغرتين.
أي جانب يحمل أي منطق؟
\end{exercise}

\begin{exercise}[$\star\star$]\label{exo:g8:reproductive-systems:9}
لماذا تكون سنوات الدورات الأولى غير منتظمة غالبًا،
ولماذا يعد اختلاف الطول أمرًا عاديًا؟ (حقيقتان من
\cref{rem:g8:reproductive-systems:living}.)
\end{exercise}

\begin{exercise}[$\star\star$]\label{exo:g8:reproductive-systems:10}
فسّر تقلصات \omterm{prop:g8:reproductive-systems:cycle}{الحيض} بمساعدة
\cref{def:g3:bones-and-muscles:muscle} --- وسمّ أمرين
يساعدان.
\end{exercise}

\begin{exercise}[$\star\star$]\label{exo:g8:reproductive-systems:11}
يجري المصنعان على جدولين متقابلين --- مستمر في مقابل
دوري بمخزون من \omterm{def:g2:born-grow-die:cycle}{الولادة}. اربط كل واحد منهما ببنية مشيجه
ودوره.
\end{exercise}

\begin{exercise}[$\star\star\star$]\label{exo:g8:reproductive-systems:12}
``الدورة هي \omterm{prop:g5:human-reproduction:plan}{الرحم} يسأل، كل شهر، السؤال الذي لا يجيب
عنه إلا \omterm{def:g5:flower-to-fruit:steps}{الإخصاب}.'' افتح الصورة على الأحداث الأربعة
المرقّمة في \cref{prop:g8:reproductive-systems:cycle}.
\end{exercise}

\section{مسألة: نماذج العيادة التعليمية}

\begin{problem}[{مسألة نهاية الأسبوع --- شرح الجهازين
بنموذجين وتقويم}]\label{pb:g8:reproductive-systems:1}
يستعمل يوم مفتوح في مركز صحي ثلاث وسائل تعليمية: نموذج
للجهاز الذكري، ونموذج للجهاز الأنثوي، وتقويم كبير لدورة
من 28 يومًا. وأنت المتطوع الشارح.

\medskip\noindent\textbf{الجزء الأول --- النموذجان.}
\begin{enumerate}
  \item يتتبع زائر النموذج الذكري من \omterm{def:g8:reproductive-systems:male}{الخصيتين} إلى المخرج.
        اسرد الطريق، وسمّ ما يضاف في أثنائه وترتيب العضو
        ذي المهمتين.
  \item على النموذج الأنثوي، بيّن أين ينتظر مخزون
        \omterm{def:g8:sexual-reproduction:gametes}{الأمشاج}، وأين يحدث \omterm{def:g5:flower-to-fruit:steps}{الإخصاب} عادة، وأين يستقر
        \omterm{def:g5:human-reproduction:pregnancy}{الحمل}.
  \item يسأل زائر لماذا يختلف مصنعا \omterm{def:g8:sexual-reproduction:gametes}{الأمشاج} في النموذجين
        هذا الاختلاف الحاد --- ملايين يوميًا في مقابل
        واحدة في الشهر. أجب بجملة الختام في
        \cref{ex:g8:reproductive-systems:egg}.
  \item ويسأل آخر ماذا كان الجهازان يفعلان قبل \omterm{def:g8:puberty:puberty}{البلوغ}.
        أجب بآلة \cref{ch:g8:puberty} في سطر واحد.
\end{enumerate}

\medskip\noindent\textbf{الجزء الثاني --- التقويم.}
\begin{enumerate}[resume]
  \item امش على التقويم: ضع \omterm{prop:g8:reproductive-systems:cycle}{الحيض}، وإعادة البناء،
        \omterm{prop:g8:reproductive-systems:cycle}{والإباضة}، واليوم الخصب، بتواريخها التقريبية.
  \item يفترض زائر أن اليوم 28 هو نفسه عند الجميع. صحّح
        الافتراض وفق \cref{rem:g8:reproductive-systems:living}
        --- متغيّران اثنان.
  \item بيّن على التقويم ما الذي يتغير إذا حدث \omterm{def:g5:flower-to-fruit:steps}{الإخصاب}
        قرب اليوم 14 --- أي الأحداث يُلغى، وأيها يبدأ
        (\cref{def:g5:human-reproduction:pregnancy} يكمل
        الحكاية).
  \item تسأل مراهقة، بهدوء، هل تعني التقلصات أن شيئًا ما
        على غير ما يرام. أعط جواب \omterm{def:g3:bones-and-muscles:muscle}{العضلة} والمساعدتين ---
        والعنوان الذي يقصد لما وراءهما.
\end{enumerate}

\medskip\noindent\textbf{الجزء الثالث --- حرفة
الشارح.}
\begin{enumerate}[resume]
  \item يضحك الزوار من النماذج. وجملتك الافتتاحية تعامل
        الجهازين كما تعاملهما الفقرة الأولى من
        \cref{ch:g8:reproductive-systems}. اكتبها.
  \item لخّص الجهازين في المقارنة ذات السطور الخمسة في
        \cref{met:g8:reproductive-systems:sides}، عن ظهر
        \omterm{def:g4:heart-and-blood:heart}{قلب}.
  \item يسأل زائر أين تقع بدايات المواليد على الوسائل.
        اربط الوسائل الثلاث في سلسلة تعطي الخطوات الثلاث
        في \cref{prop:g8:sexual-reproduction:core}.
  \item اختم اليوم المفتوح بجملة واحدة عن سبب تعليم
        المراكز الصحية هذا التشريح أصلًا --- المعرفة في
        مقابل الأسطورة، كما حاجّ \cref{ex:g8:puberty:questions}.
\end{enumerate}
\end{problem}
